% Body-facing result-table skeleton v6. One-page review sheet, not a full paper.
% Historical planning version: artifacts/aamas2027/experiment_matrices_20261008_v3/.
% Every empirical cell is intentionally blank until its confirmation card is executed.
\documentclass[sigconf,anonymous]{aamas}
\usepackage{amsmath}
\usepackage{booktabs}
\usepackage{array}
\usepackage{makecell}
\usepackage{url}
% Copyright and conference block copied from the official AAMAS 2027 sample.
\makeatletter
\gdef\@copyrightpermission{
  \begin{minipage}{0.2\columnwidth}
   \href{https://creativecommons.org/licenses/by/4.0/}{\includegraphics[width=0.90\textwidth]{by}}
  \end{minipage}\hfill
  \begin{minipage}{0.8\columnwidth}
   \href{https://creativecommons.org/licenses/by/4.0/}{This work is licensed under a Creative Commons Attribution International 4.0 License.}
  \end{minipage}
  \vspace{5pt}
}
\makeatother
\setcopyright{ifaamas}
\acmConference[AAMAS 2027]{Proc.\@ of the 26th International Conference
on Autonomous Agents and Multiagent Systems (AAMAS 2027)}{3--7 May 2027}{Hanoi, Vietnam}{M.~Baldoni, F.~Fang, W.~Yeoh, N.~Yorke-Smith (eds.)}
\copyrightyear{2027}
\acmYear{2027}
\acmDOI{}
\acmISBN{}
\acmSubmissionID{INTERNAL REVISION}
\submissionType{Research Paper Track}


\vfuzz=1.5pt
\acmSubmissionID{INTERNAL PRE-RESULTS -- NOT FOR SUBMISSION}
\submissionType{Research Paper Track}

\title[Body-Facing Experiment Tables]{Body-Facing Experiment Tables}
\author{Anonymous Authors}
\affiliation{\institution{Internal research draft}}
\email{omitted}
\ccsdesc[500]{Computing methodologies~Multi-agent systems}
\ccsdesc[300]{Computing methodologies~Cooperation and coordination}
\keywords{multi-agent systems, role formation, peer judgment, online adaptation, experiments}

\newcommand{\blank}{\textemdash}
\newcommand{\rare}{\textsc{rare}}
\newcommand{\unknown}{\textsc{unknown}}
\newcommand{\metric}[1]{\makecell[c]{#1}}
\newcommand{\dashrow}{\blank & \blank & \blank & \blank & \blank & \blank}

\begin{abstract}
This review copy fixes the compact, body-facing result-table layout. It keeps
method rows and primary endpoints visible while moving roots, seeds, budgets,
scorers, and execution metadata to the experiment manifest. Dashes reserve
empirical cells; they are not zeros or predicted values.
\end{abstract}

\begin{document}
\onecolumn
\noindent\textbf{Experiment results: layout preview}\hfill\textit{No empirical values entered}

\smallskip
\noindent\small The three grids are candidates for the body results section.
\rare{} is the candidate method; Ctx. trust denotes contextual-trust-linear.
$U$ is future quality--complete-cost utility; $\Delta A$ is assignment-change
rate. A dash is unfilled; planned entries use stream-level 95\% intervals.

\begin{table}[!ht]
\centering
\small
\setlength{\tabcolsep}{4pt}
\renewcommand{\arraystretch}{1.12}
\begin{tabular}{@{}l*{6}{>{\centering\arraybackslash}p{0.10\textwidth}}@{}}
\toprule
Method & \multicolumn{2}{c}{ArtifactRole: prediction} & \multicolumn{3}{c}{ArtifactRole: allocation} & Streams\\
\cmidrule(lr){2-3}\cmidrule(lr){4-6}
 & Brier $\downarrow$ & Calib. $\downarrow$ & $U\uparrow$ & $\Delta A\uparrow$ & Cost $\downarrow$ & $n$\\
\midrule
Uniform & \blank & \blank & \blank & \blank & \blank & \blank\\
No update & \blank & \blank & \blank & \blank & \blank & \blank\\
Raw acceptance & \blank & \blank & \blank & \blank & \blank & \blank\\
Terminal only & \blank & \blank & \blank & \blank & \blank & \blank\\
Ctx.-trust linear & \blank & \blank & \blank & \blank & \blank & \blank\\
Pooled controller & \blank & \blank & \blank & \blank & \blank & \blank\\
Meta-Team L2-public$^{\dagger}$ & \dashrow\\
\midrule
\rare{} & \blank & \blank & \blank & \blank & \blank & \blank\\
\bottomrule
\end{tabular}
\caption{ArtifactRole prediction and future allocation (RQ1--RQ2; H1--H2). $^{\dagger}$The published adapter is not yet qualified; its cells remain unfilled.}
\label{tab:body-main-results}
\end{table}

\begin{table}[!ht]
\centering
\small
\setlength{\tabcolsep}{4pt}
\renewcommand{\arraystretch}{1.12}
\begin{tabular}{@{}l*{7}{>{\centering\arraybackslash}p{0.09\textwidth}}@{}}
\toprule
Method & \multicolumn{3}{c}{PeerSelect: service} & \multicolumn{3}{c}{PeerSelect: adaptation} & Streams\\
\cmidrule(lr){2-4}\cmidrule(lr){5-7}
 & \metric{Update p95\\(ms) $\downarrow$} & Backlog $\downarrow$ & \metric{State\\(KiB) $\downarrow$} & Payoff $\uparrow$ & Forget $\downarrow$ & \metric{Recovery\\(steps) $\downarrow$} & $n$\\
\midrule
Local uniform & \dashrow & \blank\\
No history & \dashrow & \blank\\
History only & \dashrow & \blank\\
Graph-IPD reference & \dashrow & \blank\\
\midrule
\rare{} & \dashrow & \blank\\
\bottomrule
\end{tabular}
\caption{Online service and drift adaptation on PeerSelect (RQ3--RQ4; H3). Rows share the same local menu, feedback delay, arrival rate, and resource budget.}
\label{tab:body-service}
\end{table}

\begin{table}[!ht]
\centering
\small
\setlength{\tabcolsep}{4pt}
\renewcommand{\arraystretch}{1.10}
\begin{tabular}{@{}l*{4}{>{\centering\arraybackslash}p{0.135\textwidth}}@{}}
\toprule
\multicolumn{5}{l}{\textbf{(a) Mechanism ablations}}\\
Variant & $U\uparrow$ & $\Delta A\uparrow$ & \metric{False attrib.\\$\downarrow$} & Cost $\downarrow$\\
\midrule
Full \rare{} & \blank & \blank & \blank & \blank\\
No public evidence & \blank & \blank & \blank & \blank\\
No delayed credit & \blank & \blank & \blank & \blank\\
Neither evidence nor credit & \blank & \blank & \blank & \blank\\
No ownership gate & \blank & \blank & \blank & \blank\\
No locality & \blank & \blank & \blank & \blank\\
No selected-only rule & \blank & \blank & \blank & \blank\\
No correction / replay & \blank & \blank & \blank & \blank\\
\midrule
\multicolumn{5}{l}{\textbf{(b) Responsibility stress: false producer-attribution rate $\downarrow$}}\\
Variant & \metric{Recipient\\owned} & \metric{Mixed\\ownership} & Unobserved & \metric{Resource\\failure}\\
\midrule
Full \rare{} & \blank & \blank & \blank & \blank\\
No ownership gate & \blank & \blank & \blank & \blank\\
\bottomrule
\end{tabular}
\caption{Mechanism ablations and responsibility stress on ArtifactRole (RQ2, RQ4; H4). Each contrast uses paired task streams and independent outcomes.}
\label{tab:body-safety-ablation}
\end{table}

\noindent\footnotesize\textit{Scope.} ArtifactRole and PeerSelect are complementary
tracks, not two interchangeable benchmarks. Independent-root and drift breakdowns,
secondary metrics, denominator rules, and full execution cards remain in the
experiment manifest. Benchmark qualification and baseline qualification remain
open; this preview changes presentation only.

\end{document}
